% =====================================================================
%  CARNET D'ENTRAINEMENT — FICHE 6 — THÈME 3
%  Niveau FACILE — Corrigé
% =====================================================================
\documentclass[11pt,a4paper]{article}
\usepackage[utf8]{inputenc}
\usepackage[T1]{fontenc}
\usepackage{lmodern}
\usepackage[french]{babel}
\usepackage{amsmath,amssymb,mathtools}
\usepackage{icomma}
\usepackage[version=4]{mhchem}
\usepackage{siunitx}
\sisetup{output-decimal-marker={,}}
\usepackage{xcolor}
\usepackage{colortbl}
\usepackage{tikz}
\usepackage[most]{tcolorbox}
\usepackage{enumitem}
\usepackage{array}
\usepackage{booktabs}
\usepackage{fancyhdr}
\usepackage[hidelinks]{hyperref}
\usetikzlibrary{arrows.meta,positioning,calc,shapes.geometric}
\usepackage[a4paper,margin=2.1cm,headheight=26pt,headsep=8pt]{geometry}

\definecolor{primary}{RGB}{146,95,160}
\definecolor{primarydark}{RGB}{92,52,104}
\definecolor{corcol}{RGB}{0,114,178}
\definecolor{astucecol}{RGB}{198,93,9}

\tcbset{boxbase/.style={enhanced,breakable,boxrule=0.9pt,arc=2.5pt,
   left=3mm,right=3mm,top=2.2mm,bottom=2.2mm,
   fonttitle=\bfseries,coltitle=white,toptitle=0.6mm,bottomtitle=0.6mm}}
\newtcolorbox[auto counter]{cor}[1][]{boxbase,colback=corcol!5,colframe=corcol,
   title={\textsc{Corrigé}~\thetcbcounter\ifx\relax#1\relax\relax\else\textmd{ —\itshape\ #1}\fi}}
\newtcolorbox{astucebox}[1][]{boxbase,colback=astucecol!8,colframe=astucecol,
   title={\textsc{Astuce}\ifx\relax#1\relax\relax\else\textmd{ : \itshape #1}\fi}}

\pagestyle{fancy}\fancyhf{}
\fancyhead[L]{\small\textcolor{primarydark}{\textbf{Fiche 6 · Thème 3} — Corrigé}}
\fancyhead[R]{\small\textcolor{primarydark}{\textsc{Facile}}}
\fancyfoot[C]{\small\thepage}
\renewcommand{\headrulewidth}{0.8pt}
\renewcommand{\headrule}{\hbox to\headwidth{\color{primary}\leaders\hrule height \headrulewidth\hfill}}
\setlength{\parindent}{0pt}\setlength{\parskip}{3pt}

\newcommand{\rep}{\textbf{\textcolor{corcol}{Réponse : }}}

\begin{document}
\begin{center}
{\Large\bfseries\color{primarydark}Thème 3 — Corrigé du niveau Facile}
\end{center}
\medskip

\begin{cor}[appliquer la formule]
On applique $q_{\mathrm{f}} = N_v - 2\times(\text{doublets libres}) - (\text{nb liaisons})$.
\begin{enumerate}[label=\alph*),leftmargin=2em,itemsep=2pt]
  \item $q_{\mathrm{f}} = 6 - 2\times 2 - 2 = 6 - 4 - 2 = \mathbf{0}$.
  \item $q_{\mathrm{f}} = 6 - 2\times 3 - 1 = 6 - 6 - 1 = \mathbf{-1}$.
  \item $q_{\mathrm{f}} = 4 - 2\times 0 - 4 = 4 - 0 - 4 = \mathbf{0}$.
  \item $q_{\mathrm{f}} = 5 - 2\times 0 - 4 = 5 - 0 - 4 = \mathbf{+1}$.
\end{enumerate}
\rep a) $0$ ; \; b) $-1$ ; \; c) $0$ ; \; d) $+1$.
\end{cor}

\begin{cor}[la charge formelle de l'hydrogène]
\begin{enumerate}[label=\alph*),leftmargin=2em,itemsep=2pt]
  \item $q_{\mathrm{f}}(\ce{H}) = 1 - 2\times 0 - 1 = 1 - 0 - 1 = \mathbf{0}$.
  \item \rep Tout hydrogène « normal », engagé dans \textbf{une seule liaison simple} et sans
        doublet libre, a une charge formelle \textbf{nulle}. C'est pourquoi, en pratique, on ne se
        préoccupe que des atomes lourds (\ce{C}, \ce{N}, \ce{O}\ldots) : les \ce{H} ne contribuent pas.
\end{enumerate}
\end{cor}

\begin{cor}[la somme des charges formelles donne la charge]
La somme des charges formelles vaut la charge globale de l'édifice.
\begin{enumerate}[label=\alph*),leftmargin=2em,itemsep=2pt]
  \item \ce{NH4+} : $(+1) + 4\times 0 = \mathbf{+1}$ \; $\Rightarrow$ ion de charge $+1$. \checkmark
  \item \ce{CO3^2-} : $0 + 0 + (-1) + (-1) = \mathbf{-2}$ \; $\Rightarrow$ ion de charge $-2$. \checkmark
  \item \ce{OH^-} : $(-1) + 0 = \mathbf{-1}$ \; $\Rightarrow$ ion de charge $-1$. \checkmark
\end{enumerate}
\rep a) $+1$ ; \; b) $-2$ ; \; c) $-1$.
\end{cor}

\begin{cor}[la charge formelle de l'oxygène selon sa situation]
Toujours avec $N_v = 6$ :
\begin{enumerate}[label=\alph*),leftmargin=2em,itemsep=2pt]
  \item $q_{\mathrm{f}} = 6 - 2\times 2 - 2 = 6 - 4 - 2 = \mathbf{0}$ (oxygène « normal », neutre).
  \item $q_{\mathrm{f}} = 6 - 2\times 3 - 1 = 6 - 6 - 1 = \mathbf{-1}$.
  \item $q_{\mathrm{f}} = 6 - 2\times 1 - 3 = 6 - 2 - 3 = \mathbf{+1}$ (l'oxygène de \ce{H3O+}).
\end{enumerate}
\rep a) $0$ ; \; b) $-1$ ; \; c) $+1$. On retient : un \ce{O} qui \emph{gagne} une liaison
(et perd un doublet libre) voit sa charge formelle \emph{augmenter} de $+1$.
\end{cor}

\begin{cor}[vrai ou faux ?]
\begin{enumerate}[label=\alph*),leftmargin=2em,itemsep=2pt]
  \item \rep \textbf{Faux}. La charge formelle est une charge \emph{fictive} de comptabilité, obtenue
        en partageant les liaisons à parts égales ; elle ne mesure pas la charge réelle de l'atome.
  \item \rep \textbf{Vrai}. La somme des charges formelles est toujours égale à la charge globale de
        l'édifice — c'est le test de cohérence d'une structure.
  \item \rep \textbf{Vrai}. À octet respecté, la meilleure structure est celle dont les charges
        formelles sont les plus proches de zéro (charges minimales).
\end{enumerate}
\end{cor}

\begin{astucebox}[le réflexe de la charge formelle]
$q_{\mathrm{f}} = N_v - 2\times(\text{doublets non liants}) - (\text{nb de liaisons})$, et
\textbf{une double liaison compte pour 2 liaisons}. Vérifie toujours : la somme des $q_{\mathrm{f}}$
doit redonner la charge de l'édifice.
\end{astucebox}

\end{document}
